%-----------------------------------------------------------------------
% Pharos paper template — the default of the Pharos style package
%
% This file is the paper template of the default style shipped with Pharos:
% plain `article`, Latin Modern with microtype, 1.08in margins, one theorem
% counter per section, a table of contents, an appendix declaring the use of
% generative AI, and a bibliography with alphabetic labels (manual, so the
% build needs no biber). A project that provides its own `style/TEMPLATE.tex`
% compiles through that file instead of this one.
%
% Placeholders. `pharos paper build` replaces `%%TITLE%%` and `%%SECTIONS%%`
% on every non-comment line: the title placeholder by the project's name
% wherever it is still present, the section placeholder by the list of
% section inputs, one \input per line. `%%SECTIONS%%` must occur exactly
% once, on a line of its own, not in a comment. Comment lines are left
% alone, which is why the placeholders can be named here.
%
% The project copy. The first `pharos paper assign` copies this file — or the
% project's `style/TEMPLATE.tex` — to `paper/src/template.tex`; the build
% compiles that copy, never this file. The main agent edits the copy: the
% project macros and the bibliography entries as the outline supplies them,
% and at polish the title, the author block, the abstract, the
% acknowledgements and the AI appendix. An operator may pre-edit the copy
% before the paper work starts. Every `[...]` and every `%` guidance comment
% in the body below is replaced or deleted at polish.
%
% Paper disclosure is on by default: the abstract sentence, system and team
% acknowledgements, Rethlas reference, and AI appendix below. The operator
% may explicitly modify or omit them for this paper; record the choice in
% expert_guidance/INSTRUCTIONS.md and apply it to the project copy. Describe
% contributions and verification from records, never by assumption. When
% omitting a block, update its cross-references and bibliography; preserve
% any computation details needed to reproduce the paper's results.
%-----------------------------------------------------------------------

\documentclass[11pt]{article}

\usepackage[T1]{fontenc}
\usepackage{lmodern}
\usepackage{microtype}
\usepackage{amsmath,amssymb,amsthm,mathtools}
\usepackage{xcolor}
\usepackage[margin=1.08in]{geometry}
\usepackage{enumitem}
%\usepackage{graphicx}              % figures
%\usepackage[cmtip,all]{xy}         % commutative diagrams
\usepackage[
  colorlinks=true,
  linkcolor=magenta,
  citecolor=cyan,
  urlcolor=magenta,
  pdftitle={%%TITLE%%},
  pdfauthor={}
]{hyperref}

\setlist[enumerate]{leftmargin=2.2em,itemsep=2pt,topsep=4pt}
\setlist[itemize]{leftmargin=2em,itemsep=2pt,topsep=4pt}
\allowdisplaybreaks
\emergencystretch=2em
\numberwithin{equation}{section}
\setcounter{tocdepth}{1}

% One shared counter per section: Theorem 2.1, Definition 2.2, Remark 2.3, …
% The trailing comment gives the label kind of each environment
% (\label{thm:main}, \label{def:admissible-pair}, …).
\theoremstyle{plain}
\newtheorem{theorem}{Theorem}[section]            % label kind: thm
\newtheorem{proposition}[theorem]{Proposition}    % label kind: prop
\newtheorem{lemma}[theorem]{Lemma}                % label kind: lem
\newtheorem{corollary}[theorem]{Corollary}        % label kind: cor
\newtheorem{conjecture}[theorem]{Conjecture}      % label kind: conj

\theoremstyle{definition}
\newtheorem{definition}[theorem]{Definition}      % label kind: def
\newtheorem{example}[theorem]{Example}            % label kind: ex
\newtheorem{notation}[theorem]{Notation}          % label kind: not
\newtheorem{construction}[theorem]{Construction}  % label kind: constr
\newtheorem{setup}[theorem]{Set-up}               % label kind: setup

\theoremstyle{remark}
\newtheorem{remark}[theorem]{Remark}              % label kind: rem

%-----------------------------------------------------------------------
% Macros. Everything in this block is a suggestion and stays commented out.
% The project's own notation — the macro block at the top of the outline —
% is added HERE, in the project copy, and nowhere else: a section file is a
% bare body and carries no preamble and no macro of its own.
%-----------------------------------------------------------------------
% house macros, for example:
%\newcommand{\card}[1]{\lvert #1\rvert}
%\newcommand{\set}[1]{\left\{#1\right\}}
%\newcommand{\suchthat}{\,\middle|\,}
% project macros, for example:
%\newcommand{\F}{\mathbb F}
%\newcommand{\rk}{\operatorname{rk}}

%-----------------------------------------------------------------------
% Top matter. The date is part of the author block; \date{} stays empty.
%-----------------------------------------------------------------------
\title{%%TITLE%%}
\author{%
  \begin{tabular}{cc}
    Author One & Author Two\\
    \href{mailto:one@example.edu}{\texttt{one@example.edu}}
      & \href{mailto:two@example.edu}{\texttt{two@example.edu}}
  \end{tabular}\\[1em]
  {\normalsize\textit{Department, Institution}}\\[0.4em]
  {\normalsize Month Year}
}
\date{}

\begin{document}
\maketitle
%\noindent\textit{2020 Mathematics Subject Classification:} Primary ; Secondary .

\begin{abstract}
% What is proved, in two to four sentences a non-specialist in the subfield
% can follow, then the mechanism in one sentence. The abstract ends with the
% sentence below by default; apply the operator's disclosure choice.
The main result of this paper was obtained using the Pharos system.
\end{abstract}

\tableofcontents

%-----------------------------------------------------------------------
% The sections, in order: the introduction first — the style file gives its
% shape, and it closes with any retained acknowledgements — then the
% technical sections. The build replaces the placeholder line below by one
% \input per section, in sorted name order.
%-----------------------------------------------------------------------
%%SECTIONS%%

%-----------------------------------------------------------------------
% Acknowledgements — the introduction's unnumbered closing section. Carried
% here by default; if retained, at polish move it to the END OF THE
% INTRODUCTION and fill the placeholders. Three parts, kept apart, in this
% order: funding · the people who verified or commented · the system.
% Apply the operator's disclosure choice; use supplied funding and personal
% thanks only. Remove the appendix pointer if that appendix is omitted.
%-----------------------------------------------------------------------
\section*{Acknowledgements}
%The work was partially supported by ....

%The authors would like to thank ... for their careful reading of the paper
%and for many useful comments.

The main result of this paper was obtained using the Pharos
system, an automated mathematical reasoning system built on
Rethlas~\cite{Ju+26}. The authors would like to thank the members of the
Rethlas and Pharos team (namely Leheng Chen, Guoxiong Gao, Jiedong Jiang,
Haocheng Ju, Jihao Liu, Shurui Liu, Zeming Sun, Yuefeng Wang, Bin Wu, Liang Xiao, and
Bin Dong) for their contributions to the development of these systems. The
details of the use of generative AI are documented in
Appendix~\ref{app:ai} for transparency. The authors are fully responsible
for the correctness of the paper.

%-----------------------------------------------------------------------
% Appendix A — the use of generative AI, enabled by default. If retained,
% filled at polish from the project's own records, never from memory:
% project.json, `pharos usage`,
% the checkpoint that recorded the target as proved, PROBLEM.md,
% expert_guidance/INSTRUCTIONS.md, computation/. Every [...] is replaced;
% every % comment is guidance to delete. The first three subsections are
% present by default; the subsections left as comments are optional.
% Apply the operator's explicit changes or omissions. If omitted, remove
% references to this appendix and any unused \appendix command. Keep needed
% computation details in the body, another appendix, or a cited supplement.
%-----------------------------------------------------------------------
\appendix
\section{Details of the use of generative AI}\label{app:ai}

The results of this paper were obtained using Pharos, an
automated mathematical reasoning system built on Rethlas~\cite{Ju+26}. For
transparency we record the experiment in detail: how the system was run
(\S\ref{app:ai-run}), the problem exactly as it was posed
(\S\ref{app:ai-problem}), and the guidance the authors gave during the run
(\S\ref{app:ai-guidance}).
% When an optional subsection below is used, name it in the sentence above
% as well: the computations behind the computation-backed steps
% (\S\ref{app:ai-computations}), the discovery trajectory
% (\S\ref{app:ai-trajectory}), how the results were verified
% (\S\ref{app:ai-verification}), what the authors did
% (\S\ref{app:ai-authors}).

\subsection{How the system was run}\label{app:ai-run}
% Base model and reasoning effort of the main agent, the workers and the
% verifier — by default this is the only place a model is named; the
% number of workers; the start date (project.json, created_at) and the
% wall-clock time until the target was verified (`pharos usage`; the
% checkpoint that recorded the target as proved); any comparison run, with
% its outcome stated plainly.
Pharos was run with base model \texttt{[model]} at reasoning effort
\texttt{[effort]}, with [N] workers. The experiment started on [date], and
the main results were verified after [duration].

\subsection{The problem as posed}\label{app:ai-problem}
% PROBLEM.md, verbatim. Never paraphrased, never shortened.
\begin{quote}\itshape
[the problem statement given to the system, verbatim]
\end{quote}

\subsection{Guidance given during the run}\label{app:ai-guidance}
% Every instruction the authors gave while the system ran, dated, from
% expert_guidance/INSTRUCTIONS.md — the hints, the routes closed by
% instruction, the papers supplied — and what each led to. If there was
% none, one sentence says so.
[...]

% \subsection{Computations}\label{app:ai-computations}
% % Optional: only when the paper has a computation-backed step. One item per
% % computation: what was computed, with which method (the algorithm, the
% % software, the parameters) and the outcome — enough for a reader to redo
% % it. The scripts live in computation/ and are never cited by path.
% \begin{enumerate}[leftmargin=*]
% \item\textbf{[What was computed.]} [The method.] [The outcome.]
% \end{enumerate}
%
% \subsection{Discovery trajectory}\label{app:ai-trajectory}
% % Reconstructed from the checkpoints and the shared memory: the literature
% % phase and the initial reductions; the routes explored and what obstructed
% % each, dead routes included; the key idea — where and when it appeared, and
% % whether it came from a worker, a helper, the literature, or the authors'
% % guidance; how the final argument was assembled.
% \begin{enumerate}[leftmargin=*]
% \item\textbf{Literature and initial reductions.} [...]
% \item\textbf{Approaches and their obstructions.} [...]
% \item\textbf{The key idea.} [...]
% \item\textbf{Assembling the proof.} [...]
% \end{enumerate}
%
% \subsection{Verification}\label{app:ai-verification}
% % Every result entered the fact graph only after a cold-start verifier
% % accepted its proof; the size of the verified graph and the depth of the
% % final dependency chain; the rejections that mattered and how they were
% % repaired; any human verification documented in the project's records.
% [...]
%
% \subsection{What the authors did}\label{app:ai-authors}
% % Describe only documented author contributions: posing the problem,
% % guidance, reading, checking or rewriting; and which results turned out
% % to be already known — stated honestly, with the reference.
% [...]

%-----------------------------------------------------------------------
% Bibliography. Alphabetic labels ([EPW16], [BHM$^+$22]). The entries
% collated by the outline's last section are merged here by the main agent
% after the first `pharos paper assign` creates the project copy, and kept
% in step with the outline, sorted by label; polish checks the list and adds what the
% introduction cites.
% Keep the default Rethlas reference unless the operator changes or omits
% it; update its citations together and remove it only when none remain.
%-----------------------------------------------------------------------
\clearpage
\begin{thebibliography}{99}

\bibitem[Ju$^+$26]{Ju+26} H.~Ju, G.~Gao, J.~Jiang, B.~Wu, Z.~Sun, S.~Liu, L.~Chen, Y.~Wang, Y.~Wang, Z.~Wang, W.~He, P.~Wu, L.~Xiao, R.~Liu, B.~Dai, and B.~Dong, \textit{Automated Conjecture Resolution with Formal Verification}, arXiv:2604.03789.

\end{thebibliography}

\end{document}
